% Part: incompleteness
% Chapter: representability-in-q
% Section: c

\documentclass[../../../include/open-logic-section]{subfiles}

\begin{document}

\olfileid{inc}{req}{cee}
\olsection{$C$ ప్రమేయాలు}

క్రింది ప్రమేయాలను కలిగి ఉన్న అతి చిన్న ప్రమేయాల సమితి $C$ అనుకుందాం:
\begin{enumerate}
\item $0$,
\item ఉత్తరవర్తి ప్రమేయం,
\item కూడిక,
\item గుణకారం,
\item ప్రక్షేపణలు, మరియు
\item సమానత్వానికి లక్షణ ప్రమేయం, $\Char{=}$;
\end{enumerate}
అలాగే ఈ క్రింది క్రియల కింద సంవృతమైనది:
\begin{enumerate}
\item సంయోజనం, మరియు
\item క్రమబద్ధ ప్రమేయాలకు వర్తించే అపరిమిత శోధన.
\end{enumerate}
చివరి పరిమితి అర్థం, ఫలితం మొత్తం ప్రమేయమైతేనే $\mu$ క్రియను వాడవచ్చనేది. దీన్ని \emph{సాధారణ పునరావృత} ప్రమేయాల నిర్వచనంతో పోల్చండి: ఇక్కడ కూడిక, గుణకారం, $\Char{=}$లను చేర్చాం; కానీ ఆదిమ పునరావృతాన్ని తొలగించాం.

$C$లోని ప్రతిదీ పునరావృతమేనని స్పష్టమే; ఎందుకంటే కూడిక, గుణకారం, $\Char{=}$ పునరావృతాలు. దీనికి విలోమం కూడా నిజమని చూపుతాం. అంటే $C$లోని ఇతర అంశాలతో ఆదిమ పునరావృతాన్ని నిర్వహించగలమన్న మాట.

\end{document}
